\section{本章小结}
\label{sec:rvd-summary}

本章从“由试验结果算出的数”出发。随机变量是样本空间上的函数，它把结果上的概率搬到实轴上，得到分布。分布只记录取值落在各处的概率，不记录这些取值来自哪个结果，所以同分布的变量可以在每次试验中取值完全不同，涉及几个变量之间关系的问题要留给随机向量。

分布有几种表示，各有成立的条件。分布函数$F_X(x)=\mathbb{P}(X\leq x)$对一切随机变量都有定义，单调不减、右连续、两端趋于$0$和$1$，跳跃高度是单点概率。概率质量函数只适用于离散变量，用求和计算概率；密度只适用于有密度的变量，用积分计算概率，它的值是单位长度上的概率，可以大于$1$，在有限个点处改动不影响分布。分位数$Q(p)=\inf\{x:F(x)\geq p\}$从反方向读分布函数，对一切分布都存在，并且与分布函数相互确定。

用可能不用等的等待时间$W$把这些表示串起来，可以看清各种表示分别用在哪里。分布函数由式~\eqref{eq:rvd-mixed-cdf}给出，任何区间的概率都能从它和它的左极限读出，例如$\mathbb{P}(W\leq1)\approx0.7425$。$W$在$0$处有$0.3$的质量，这部分只能用概率表描述；其余$0.7$的质量有密度$0.7\,\mathrm{e}^{-t}$，只能用积分描述；两者都不能单独代表整个分布。分位数在$p\leq0.3$时为$0$，中位数是$\ln(7/5)\approx0.336$分钟，它们描述的是质量按比例分开的位置。

表~\ref{tab:rvd-summary}汇总了各种表示。选用哪一种，先看分布属于哪一类，再看问题是给定范围求概率，还是给定概率求阈值。

\begin{BookInlineTable}
  \centering
  \begin{tabular}{lll}
    \toprule
    表示 & 适用范围 & 计算方式\\
    \midrule
    分布函数$F_X$ & 一切分布 & 相减，端点处用左极限\\
    质量函数$p_X$ & 离散分布 & 对取值求和\\
    密度$f_X$ & 有密度的分布 & 在区间上积分\\
    跳跃与密度组合 & 如$W$、$\min\{T,2\}$ & 跳跃求和，连续部分积分\\
    分位数$Q$ & 一切分布 & 给定概率反求阈值\\
    \bottomrule
  \end{tabular}
  \caption{分布的几种表示及其适用范围。Cantor分布的分布函数连续，却既无概率表也无密度，只能用分布函数或分位数函数描述。}
  \label{tab:rvd-summary}
\end{BookInlineTable}

这些表示的边界也要记住。分布函数连续不保证密度存在；一般情形下哪些集合可以赋予概率、分布函数为什么唯一确定分布，都在第~\ref{chap:probability-spaces}章统一建立。完整的分布包含全部概率规律，但使用起来往往太细。下一章从分布中提取平均位置和波动大小这样的少数数值，并说明它们存在的条件；中位数与加权平均的对比已经表明，几个数值只是分布的概括，不能代替分布本身。
